\chapter{解题感悟}
\label{cp:notes}

第三周的实验围绕 Python 打包发布、智能体协作和机器学习训练展开，相比前两周的 Shell 命令和调试，这周的实验明显更接近真实的工程开发流程，也更有挑战。

第 9 题让我第一次完整走了一遍"源码 -- 构建 -- 安装 -- 运行"的打包流程。一开始对 pyproject.toml 里的 \verb|[build-system]|、\verb|[project]|、\verb|[project.scripts]| 这些配置完全不熟悉，不知道 \verb|sdt-greet| 这个命令是从哪里来的。后来才明白 \verb|[project.scripts]| 就是把包里的函数注册成命令行入口，构建之后 pip 安装，系统就能直接调用。第 9 题要求"只从 wheel 安装"，我一开始想在干净环境里直接 \verb|pip install greetlab|，结果当然装不上，后来才理解题目的用意：只有把构建产物本身装进去，才能证明打包是正确的。

第 10 题是让我感触最深的一题。题目要求让编程智能体进入"可验证的修复循环"：先写测试确认失败，再让智能体修复，最后人工检查 diff。我一开始是让智能体直接改代码，改完就跑，结果代码虽然能运行，但我自己都不知道改了什么、为什么改。按照题目的流程做下来，先看到 \verb|TypeError| 的失败输出，再看到智能体补上 \verb|argv| 参数和空白校验，最后用 diff 确认只有 cli.py 有改动，整个过程每一步都有证据，非常踏实。我也第一次意识到，测试不是"多写的作业"，而是让修改可验证的锚点；ai\_log.md 用几行字把提示、改动和验证记录下来的习惯也值得坚持。

第 11 题字数限制 400 字看起来不难，真正写的时候才发现把问题讲清楚还要控制字数并不容易。我一开始写 Issue 的时候把复现步骤写得过于冗长，又漏掉了环境信息，初版写完是 403 字还超了。精简的时候学会了把"环境、复现、期望、实际"这类信息用最少的词说清楚，把未知信息诚实地标注成"待确认"而不是编一个结论。这让我意识到，协作材料最重要的是让接收方"能直接执行"，而不是堆字数。

第 12 题是第一次用 PyTorch 写训练循环。题目里把 \verb|zero_grad|、\verb|backward|、\verb|step| 都留成了 TODO，补全本身不难，难的是理解为什么要 \verb|zero_grad|：不清空梯度的话，多次 backward 会把梯度累加，参数更新就会出错。运行出 \verb|final loss: 2.39e-12| 的时候很有成就感，也体会到机器学习调试"看指标、找原因、改逻辑"的基本思路。

这周的实验比前两周更"工程化"：打包、测试、协作、训练，每一步都有明确的产物和验证标准。虽然过程依然有卡壳的时候（比如环境变量没生效、numpy 缺失的警告），但按照"先验证、再修改、最后核对"的流程，问题都能定位到具体环节。希望在后面的课程里能把这种工作方式保持下去。

